% Figure: photodiode responsivity against wavelength. The ideal responsivity
% R = eta q lambda / (h c) rises linearly with wavelength at fixed quantum
% efficiency (the dashed ideal line for eta = 1), while a real detector tracks
% that line over its sensitive band and falls to zero at the band edges: below the
% short-wavelength cut where surface absorption robs the junction, and above the
% long-wavelength cut set by the semiconductor bandgap.
\begin{figure}[htbp]
\centering
\begin{tikzpicture}
\begin{axis}[
  width=10.5cm, height=6.4cm,
  xlabel={wavelength $\lambda$ (nm)},
  ylabel={responsivity $\mathcal{R}$ (A/W)},
  xmin=300, xmax=1200,
  ymin=0, ymax=1.0,
  domain=300:1200, samples=120,
  axis lines=left,
  xtick={400,600,800,1000,1200},
  ytick={0,0.2,0.4,0.6,0.8,1.0},
  legend style={font=\scriptsize, at={(0.03,0.97)}, anchor=north west, draw=none},
  tick label style={font=\scriptsize},
  label style={font=\small},
  every axis plot/.append style={very thick},
]
% ideal responsivity for eta = 1: R = q lambda /(h c) = lambda[nm]/1240 A/W
\addplot[engamber, dashed]{x/1240};
\addlegendentry{ideal, $\eta=1$}
% realistic silicon detector: eta ~ 0.85 in band, rolling off at both edges.
% model eta(lambda) as a smooth window; R = eta * lambda/1240.
\addplot[engblue]{ (0.85/(1 + exp(-(x-420)/25)) ) * (1/(1+exp((x-1050)/25))) * x/1240 };
\addlegendentry{silicon (realistic)}
% mark the peak-responsivity region
\node[engred, font=\scriptsize, anchor=south] at (axis cs:900,0.62) {peak near band edge};
\draw[engred, ->] (axis cs:900,0.60) -- (axis cs:970,0.52);
\end{axis}
\end{tikzpicture}
\caption[Photodiode responsivity against wavelength]{The responsivity of a
photodiode, the photocurrent per watt of incident optical power, is
$\mathcal{R}=\eta q\lambda/(hc)$, which for unit quantum efficiency rises linearly
with wavelength because each photon carries less energy at longer wavelength and a
fixed power delivers more photons. The dashed line is that ideal at $\eta=1$; the
solid curve is a realistic silicon detector, which tracks the ideal across its
sensitive band at a quantum efficiency near $0.85$ and falls to zero at both
edges. Below about $400\,\si{\nano\meter}$ the carriers are generated so close to
the surface that they recombine before collection; above about
$1100\,\si{\nano\meter}$ the photon energy drops below the silicon bandgap and the
material is transparent. The peak responsivity sits just short of the
long-wavelength cut, near $900\,\si{\nano\meter}$, which is why silicon detectors
pair naturally with the $905\,\si{\nano\meter}$ LIDAR source of the active-sensing
case study, and why reaching $1.55\,\si{\micro\meter}$ needs a
narrower-gap material such as indium gallium arsenide.}
\label{fig:vol04ch11:photodiode-responsivity}
\end{figure}
